% =====================================================================
% Updated sections for the MATH 629 report (report_v3).
% Each block below REPLACES the section of the same number in report_v3.
% Sections 1, 2 and the appendices' literature content are unchanged.
% Section 7 is in section7_results.tex.
% Preamble needs: \usepackage{booktabs, graphicx, amsmath}
% \cite keys (cont2014, kyle1985, gould2013, leeready1991, amihud2002) are
% placeholders: replace them with the keys in your .bib (refs [4],[13],[7],[14],[2]).
% =====================================================================


% =====================================================================
% REPLACES: Abstract
% =====================================================================
\begin{abstract}
We study whether deep learning models can extract economically meaningful
short-horizon signals from high-frequency cryptocurrency limit order book (LOB)
data, and whether PCA compression and cross-asset features help. Using 1,030,551
one-second Bitcoin--Ethereum LOB snapshots from 7--19 April 2021, we build a single
feature pipeline (raw or PCA-compressed LOB depth, cross-asset market state and six
microstructure factors). On it we evaluate a Residual CNN, a bidirectional LSTM, a
Transformer and a LightGBM baseline under one protocol: a purged
train/validation/test split, trading rules chosen on validation only, and a
non-overlapping backtest with a 1.5\,bp fee.

All models show a weak but consistent directional signal (directional F1
0.44--0.47), although their accuracy is no better than always predicting the majority
Flat class. Trading every prediction loses 39--49\%, and confidence thresholds
selected out of sample leave almost no trades. A hysteresis rule with separate entry
and exit levels makes the signal profitable under the backtest's assumptions. Over six
walk-forward test days, a seed-averaged CNN earns $+3.1\%$ to $+10.0\%$ per day
against a buy-and-hold mean of $-1.6\%$. However, this profit disappears when
execution is delayed by one to three seconds. The signal is real but decays within about
a second, so we find no tradable edge at one-second resolution after costs. Longer
prediction horizons (30--300\,s) do not help, and the deep architectures add little
directional skill over the LightGBM baseline.
\end{abstract}


% =====================================================================
% REPLACES: Section 3 Research Questions
% (RQ1-RQ4 unchanged in substance; RQ4 made precise, RQ5 added)
% =====================================================================
\section{Research Questions}
\label{sec:rq}
This project addresses the following questions:
\begin{description}
  \item[RQ1.] Can deep learning architectures (CNN, LSTM, Transformer) extract
  directional signals from one-second cryptocurrency LOB data at a 10-second horizon,
  and do they outperform a non-sequential gradient-boosting baseline?
  \item[RQ2.] Does PCA-based compression of raw LOB features improve prediction and
  trading performance compared with using the raw order-book features directly?
  \item[RQ3.] Do market state variables (rolling volatility, momentum, cross-asset ETH
  signals) and classic microstructure factors improve short-horizon prediction?
  \item[RQ4.] Can the models generate returns after a 1.5\,bp transaction cost when every
  trading parameter is chosen out of sample, and does any such return survive realistic
  execution delays?
  \item[RQ5.] How do the prediction horizon (10--300\,s) and the trading rule (argmax,
  confidence threshold, hysteresis) affect economic performance?
\end{description}


% =====================================================================
% REPLACES: Section 4 Data (4.1-4.3)
% =====================================================================
\section{Data}
\label{sec:data}

\subsection{Source and Coverage}
We use the \emph{High-Frequency Crypto Limit Order Book Data} dataset on
Kaggle\footnote{\url{https://www.kaggle.com/datasets/martinsn/high-frequency-crypto-limit-order-book-data}}, which provides
one-second LOB snapshots for BTC, ETH and ADA. We use the full BTC and ETH files.
Timestamps are floored to the second (keeping the last snapshot within a second), and
the two assets are inner-joined on that second. After dropping rows whose rolling
features are not yet defined, 1,030,551 aligned snapshots remain, covering
7 April 2021 11:34 UTC to 19 April 2021 09:54 UTC (about 12 days).

Each BTC snapshot records 15 levels on each side: bid distances $\{d^b_i\}_{i=0}^{14}$,
ask distances $\{d^a_i\}_{i=0}^{14}$, bid notionals $\{v^b_i\}_{i=0}^{14}$ and ask notionals
$\{v^a_i\}_{i=0}^{14}$, i.e.\ 60 LOB features. Here \emph{distance} is the relative offset of
level $i$ from the mid-price and \emph{notional} is the dollar value resting at that level.
From ETH we use the mid-price, spread and level-0 notionals.

\subsection{Label Construction}
For a horizon of $h$ seconds the target is the forward mid-price return
\begin{equation}
  r^{(h)}_t = \frac{p_{t+h} - p_t}{p_t},
  \label{eq:fwdret}
\end{equation}
where $p_{t+h}$ is looked up \emph{by timestamp}, not by row offset. The feed has gaps, so
``$h$ rows later'' is not always ``$h$ seconds later''. Rows whose future timestamp is
missing are never used as targets. Labels are three-class:
\begin{equation}
  y_t =
  \begin{cases}
    \text{Up}   & \text{if } r^{(h)}_t > +\delta_h,\\
    \text{Down} & \text{if } r^{(h)}_t < -\delta_h,\\
    \text{Flat} & \text{otherwise.}
  \end{cases}
  \label{eq:label}
\end{equation}
The main horizon is $h=10$\,s with $\delta_{10} = 1.5\times10^{-4}$, which gives about
21\% Up, 57\% Flat and 21\% Down. For the horizon study (Section~\ref{sec:horizon}) we also
build labels for $h \in \{30, 60, 300\}$\,s with $\delta_h = 3\times10^{-4}$, the round-trip
cost, so that an Up label marks a move that is still profitable after paying to enter and
exit. The share of Flat labels falls with the horizon (43.1\%, 35.9\%, 24.4\% and 10.7\% of
the test set at 10, 30, 60 and 300\,s).

\subsection{Train--Validation--Test Split}
\label{sec:split}
We use a strictly chronological 70/15/15 split: training (7 April 11:34 to about
15 April 20:00), validation (to about 17 April 15:00) and test (to 19 April 09:54). The
validation split is used for early stopping and for choosing every trading parameter.
The test split is evaluated once.

Instead of a fixed embargo of 10 rows, we \emph{purge} by time. A row is a valid target
only if its whole label window $[t, t+h]$ lies inside its own split. No training label is
then computed from validation or test prices, and no validation label from test prices,
for any horizon $h$. Input windows may reach back across a split boundary. This is
harmless, because all features are backward-looking.

To check that results do not depend on this single test window, Section~\ref{sec:walkforward}
also uses \emph{walk-forward} folds. The data are cut into 24-hour blocks from the first
timestamp. Fold $k$ trains on blocks $k,\dots,k+4$, validates on block $k+5$ and tests on
block $k+6$, giving six disjoint test days (13--18 April). Each fold has about 432k
training, 86k validation and 86k test targets.


% =====================================================================
% REPLACES: Section 5.1 Feature Engineering Pipeline (5.1.1-5.1.6)
% (5.1.1 and 5.1.2 are unchanged in substance and kept short here)
% =====================================================================
\subsection{Feature Engineering Pipeline}
\label{sec:features}
A single pipeline produces one feature table that all models read, so every model sees
identical features, labels and splits.

\subsubsection{Stage 1: Raw LOB Features}
The 60 BTC LOB features (distances and notionals at 15 levels per side) form the base
input and preserve the full depth profile of the book.

\subsubsection{Stage 2: Market State Features}
Fifteen features capture BTC and ETH short-term dynamics and their interaction:
\begin{itemize}
  \item BTC: spread, level-0 order imbalance, rolling volatility and momentum of log
  returns (10\,s and 60\,s windows);
  \item ETH: the same six features from the ETH book;
  \item Cross-asset: BTC--ETH momentum difference (10\,s, 60\,s) and BTC/ETH volatility
  ratio (10\,s).
\end{itemize}

\subsubsection{Stage 3: Classic Microstructure Factors}
Six factors from the microstructure literature yield 13 columns:
\begin{enumerate}
  \item \textbf{Queue imbalance} \cite{cont2014}: exponentially weighted imbalance over
  five levels, $\mathrm{QI}_t = \sum_{i=0}^{4} w_i \frac{v^b_i - v^a_i}{v^b_i + v^a_i}$ with
  $w_i \propto e^{-0.5 i}$, $\sum_i w_i = 1$.
  \item \textbf{Order flow imbalance (OFI)} \cite{cont2014}:
  $\mathrm{OFI}_t = e^b_t - e^a_t$ with
  $e^b_t = v^b_{0,t}\mathbb{1}[P^b_t \ge P^b_{t-1}] - v^b_{0,t-1}\mathbb{1}[P^b_t \le P^b_{t-1}]$
  and $e^a_t$ defined symmetrically, plus 10\,s and 60\,s running sums. Because the data
  store distances rather than prices, the best bid and ask prices are rebuilt as
  $P^b_t = p_t(1 + d^b_{0,t})$ and $P^a_t = p_t(1 + d^a_{0,t})$ and rounded to the
  \$0.01 tick before comparison. Comparing distances directly would misclassify every
  mid-price move as a queue change.
  \item \textbf{Kyle's $\lambda$} \cite{kyle1985}: rolling 60\,s regression slope of the log
  return on OFI, $\widehat{\mathrm{Cov}}(r,\mathrm{OFI}) / \widehat{\mathrm{Var}}(\mathrm{OFI})$.
  \item \textbf{Depth ratio} \cite{gould2013}: ask-to-bid notional ratio over the top
  10 levels, the corresponding imbalance, and a 10\,s rolling mean of the ratio.
  \item \textbf{Tick sign} \cite{leeready1991}: sign of the most recent non-zero mid-price
  return, forward-filled, plus 10\,s and 60\,s running sums.
  \item \textbf{Amihud illiquidity} \cite{amihud2002}: $|r_t|$ divided by the top-5-level
  book notional (book depth stands in for traded volume), plus its 60\,s rolling mean.
\end{enumerate}

\subsubsection{Stage 4: Normalization}
Every feature is standardised with a backward-looking rolling $z$-score over the last
3,600 seconds (minimum 60 observations),
\begin{equation}
  \tilde{x}_{t,j} = \operatorname{clip}\!\left(\frac{x_{t,j} - \hat\mu_{t,j}}{\hat\sigma_{t,j}},\; -10,\; 10\right),
  \label{eq:zscore}
\end{equation}
where $\hat\mu_{t,j}$ and $\hat\sigma_{t,j}$ use observations $t-3599,\dots,t$ only, so no
future information enters. Two details matter. First, no additive constant is used in the
denominator, and windows with $\hat\sigma_{t,j}=0$ map to 0. Amihud and Kyle's $\lambda$ are
$\mathcal{O}(10^{-10})$, and an $\varepsilon = 10^{-8}$ (as in an earlier version) squashes
them to zero. Second, values are clipped at $\pm 10$ because near-constant columns, such
as a spread pinned at one tick, otherwise produce extreme values whenever they move.

\subsubsection{Stage 5: PCA Compression}
A \texttt{StandardScaler} and PCA are fitted on the 60 normalized LOB features of the
\emph{training rows only} and applied to all rows. About 30 components are needed for
90\% of the variance. We keep the components with eigenvalue above one (Kaiser
criterion), capped at 10, which retain about 61\% of the variance. In the walk-forward
evaluation, the scaler and PCA are refitted on each fold's training days.

\subsubsection{Feature Sets}
Three feature sets are used throughout, with the same definition for every model
(Table~\ref{tab:featuresets}).

\begin{table}[htbp]
\centering
\caption{Feature sets. All are columns of one feature table with 1,030,551 rows.}
\label{tab:featuresets}
\begin{tabular}{lll r}
\toprule
Name & Short name & Feature groups & Features \\
\midrule
Full      & \texttt{full}      & PCA(10) + Market(15) + Classic(13) & 38 \\
No PCA    & \texttt{no\_pca}    & Raw LOB(60) + Market(15) + Classic(13) & 88 \\
No Market & \texttt{no\_market} & PCA(10) + Classic(13) & 23 \\
\bottomrule
\end{tabular}
\end{table}


% =====================================================================
% REPLACES: the "Key hyperparameters" paragraphs of 5.2.1-5.2.3
% and ADDS 5.2.4. The architecture bullet lists of 5.2.1-5.2.3 are
% unchanged and can be kept as they are.
% =====================================================================

% --- in 5.2 (intro), replace the loss paragraph with: ---
All models output a probability distribution
$\hat{\mathbf p}_t = (\hat p_{\downarrow}, \hat p_0, \hat p_{\uparrow})$ and are trained with
class-weighted cross-entropy,
\begin{equation}
  \mathcal{L} = -\sum_{c} w_c\, \mathbb{1}[y=c] \log \hat p_c, \qquad
  w_c = \frac{N}{3 N_c},
\end{equation}
where $N_c$ is the number of training targets in class $c$. All neural networks use
AdamW, gradient clipping at norm 1.0, shuffled mini-batches, and \emph{early stopping on
validation loss}. The weights of the best validation epoch are restored.

% --- 5.2.1 Residual CNN: replace the "Key hyperparameters" paragraph with: ---
\emph{Key hyperparameters:} 64 filters, $T=10$, dropout 0.3, learning rate $10^{-4}$, weight
decay $10^{-2}$, cosine annealing over at most 50 epochs, patience 10, batch size 256.
Total parameters: 118,851 for $F=38$.

% --- 5.2.2 Bidirectional LSTM: replace the "Key hyperparameters" paragraph with: ---
\emph{Key hyperparameters:} $T=10$, hidden size 32 per direction, 2 layers, dropout 0.4,
learning rate $10^{-4}$, weight decay $2\times10^{-3}$, cosine annealing over at most
40 epochs, patience 6, batch size 256.

% --- 5.2.3 Transformer: replace the "Key hyperparameters" paragraph with: ---
\emph{Key hyperparameters:} $T=30$, $d_{\text{model}}=64$, 4 heads, 3 pre-norm layers,
feed-forward width 128, dropout 0.1, OneCycleLR with maximum learning rate $10^{-3}$
(10\% warm-up, cosine decay), weight decay $10^{-4}$, at most 30 epochs, patience 6,
batch size 512. Total parameters: 105,219 for $F=38$.

% --- NEW subsection after 5.2.3: ---
\subsubsection{LightGBM Baseline}
To measure what the sequence architectures add, we train a gradient-boosted tree
classifier (LightGBM) on the features of the \emph{current second only}, with no look-back
window. It uses up to 2,000 trees, learning rate 0.05, 31 leaves, a minimum of 200
samples per leaf, 80\% row and column subsampling, balanced class weights, and early
stopping on validation multi-class log-loss (50 rounds).


% =====================================================================
% REPLACES: Section 5.3 Backtesting Framework (5.3.1-5.3.4)
% =====================================================================
\subsection{Backtesting Framework}
\label{sec:backtest}

\subsubsection{Positions, Costs and Returns}
A trading rule maps each prediction to a position $s_t \in \{-1, 0, +1\}$ (short, flat,
long), using the full account value without leverage. The position is held for exactly $h$
seconds and the net return of the step is
\begin{equation}
  r^{\text{net}}_t = s_t\, r^{(h)}_t - c\,|s_t - s_{t-1}|,
  \label{eq:netret}
\end{equation}
with $c = 1.5$\,bp per unit of position change. A reversal from $-1$ to $+1$ therefore pays
$2c$, and the first entry is also charged. Wealth compounds as
$W_T = \prod_t (1 + r^{\text{net}}_t)$. Execution is at the mid-price of the snapshot
from which the features were computed. Section~\ref{sec:delay} tests this assumption by
delaying execution.

\subsubsection{Non-Overlapping Steps}
Predictions are made every second, but each label covers $h$ seconds. Using every
prediction would count each price move up to $h$ times. We therefore keep a subset of
target rows such that each is at least $h$ seconds after the previous one. The positions
are taken at these rows and earn non-overlapping returns. For $h=10$ the test split has
15,457 such steps. Sharpe ratios are annualised with
$\sqrt{365 \times 24 \times 3600 / h}$ ($\approx 1776$ at $h=10$\,s). These values are
only comparable across models in this study, not with daily-frequency strategies.

\subsubsection{Trading Rules}
Let $\hat p_{\downarrow}$ and $\hat p_{\uparrow}$ be the predicted probabilities of Down and Up.
\begin{itemize}
  \item \textbf{Argmax}: $s_t = \arg\max_c \hat p_c - 1$, i.e.\ trade every prediction.
  \item \textbf{Threshold}: $s_t = +1$ if $\hat p_{\uparrow} > \theta$, $s_t = -1$ if
  $\hat p_{\downarrow} > \theta$, otherwise $0$.
  \item \textbf{Hysteresis}: open a long (short) position when $\hat p_{\uparrow}$
  ($\hat p_{\downarrow}$) exceeds an \emph{enter} level $\theta_{\text{in}}$ and is the larger of
  the two directional probabilities. Keep the position while its probability stays at or
  above an \emph{exit} level $\theta_{\text{out}} < \theta_{\text{in}}$, and reverse directly
  when the opposite probability exceeds $\theta_{\text{in}}$. The lower exit level avoids
  paying to close and reopen a position on every small dip in confidence.
\end{itemize}

\subsubsection{Out-of-Sample Parameter Selection}
$\theta$ is swept over 23 values in $[0.35, 0.90]$, and
$(\theta_{\text{in}}, \theta_{\text{out}})$ over $\theta_{\text{in}} \in \{0.40, 0.45,
\dots, 0.90\}$ and $\theta_{\text{out}} \in \{0, 0.05, \dots, \theta_{\text{in}} - 0.05\}$.
Both sweeps are run on the \emph{validation} split, and the configuration with the
highest validation Sharpe among those with at least 20 position changes is applied once to
the test split. An earlier version of this study selected $\theta$ on the test split, which
inflates test performance.

\subsubsection{Performance Metrics}
\emph{Statistical:} accuracy, per-class precision, recall and F1, and directional F1,
\begin{equation}
  \mathrm{DirF1} = \tfrac12\left(\mathrm{F1}_{\text{Down}} + \mathrm{F1}_{\text{Up}}\right),
\end{equation}
which ignores the dominant Flat class. Accuracy is reported next to the accuracy of always
predicting Flat.
\emph{Economic:} cumulative net return, annualised Sharpe ratio, maximum drawdown, win rate
on active steps, and turnover $\sum_t |s_t - s_{t-1}|$.
\emph{Robustness:} the test return over neighbouring threshold values, the split of P\&L
into long and short positions and by calendar day, results across walk-forward folds and
seeds, and the test return when execution is delayed by $d \in \{1,2,3,5\}$ seconds.


% =====================================================================
% REPLACES: Section 6 Experimental Setup (text + Table 2)
% =====================================================================
\section{Experimental Setup}
\label{sec:setup}
Neural networks are trained on Google Colab with an NVIDIA Tesla T4 GPU, and LightGBM on
the Colab CPU. All code (feature pipeline, models, backtest and the scripts that produce
every table) is shared across models, so differences in results come from the models, not
from preprocessing or evaluation.

Experiments on the main split use seed 42. They cover all four models and three feature
sets at $h=10$\,s, and the Full feature set at $h\in\{30,60,300\}$\,s (LightGBM: all
feature sets). The walk-forward evaluation trains the CNN with seeds 42, 43 and 44 plus
a seed-averaged ensemble, and LightGBM with seed 42, on six folds at $h=10$\,s.
Development was first done on the first $\approx$1.2 days of data (100,000 rows) and then
scaled to the full dataset. That development subset lies entirely inside the main
training split.

\begin{table}[htbp]
\centering
\caption{Hyperparameters by model. All neural networks use AdamW, gradient clipping at
1.0, shuffled batches, class-weighted cross-entropy and early stopping on validation
loss.}
\label{tab:hparams}
\begin{tabular}{lcccc}
\toprule
 & CNN & BiLSTM & Transformer & LightGBM \\
\midrule
Input window $T$ & 10 & 10 & 30 & 1 \\
Width            & 64 / 128 filters & $32\times2$ & $d_{\text{model}}=64$ & 31 leaves \\
Depth            & 2 residual blocks & 2 layers & 3 layers & $\le$ 2,000 trees \\
Dropout          & 0.3 & 0.4 & 0.1 & -- \\
Learning rate    & $10^{-4}$ & $10^{-4}$ & $10^{-3}$ (max) & 0.05 \\
Schedule         & cosine & cosine & OneCycle & -- \\
Weight decay     & $10^{-2}$ & $2\times10^{-3}$ & $10^{-4}$ & $\lambda_2 = 1$ \\
Batch size       & 256 & 256 & 512 & -- \\
Max epochs       & 50 & 40 & 30 & -- \\
Early stopping   & val loss, 10 & val loss, 6 & val loss, 6 & val log-loss, 50 \\
Parameters ($F=38$) & 118,851 & 45,827 & 105,219 & -- \\
\bottomrule
\end{tabular}
\end{table}


% =====================================================================
% REPLACES: Section 8 Conclusion
% =====================================================================
\section{Conclusion}
\label{sec:conclusion}
We built a single, leakage-controlled pipeline for one-second BTC limit order book
prediction and evaluated three deep architectures and a gradient-boosting baseline under
the same out-of-sample protocol. Our answers to the research questions are as follows.

\begin{enumerate}
  \item \textbf{Directional signal (RQ1).} All models learn a weak but consistent
  directional signal: directional F1 is 0.44--0.47 on the main split and 0.41--0.47 across
  six walk-forward folds. Accuracy (41--44\%) is no better than always predicting Flat.
  The Transformer has the highest directional F1 (0.466). However, a LightGBM model
  that sees only the current second matches it (0.465), so the sequence architectures add
  little beyond the hand-crafted features.
  \item \textbf{PCA and market features (RQ2, RQ3).} Replacing the 60 raw LOB columns by
  10 principal components changes directional F1 by less than 0.007. Removing the
  market-state block lowers it for every model (by up to 0.020). Cross-asset information
  helps slightly, while PCA compression neither helps nor hurts.
  \item \textbf{Profitability after costs (RQ4).} Trading every prediction loses 39--49\%
  at 1.5\,bp. Confidence thresholds chosen on validation leave almost no trades. In
  particular, the $+2.7\%$ CNN result of our earlier draft does not survive out-of-sample
  threshold selection. A hysteresis rule makes the signal profitable under instantaneous
  execution. The CNN earns $+26.8\%$ on the main test window, and a 3-seed CNN ensemble is
  positive on all six walk-forward test days (mean $+7.0\%$ per day against $-1.6\%$ for
  buy-and-hold). This profit, however, is concentrated in the first second after each
  snapshot. With a one-second execution delay the ensemble's mean daily return falls to
  $+0.9\%$, and by three seconds every model's return is negative. We also found that ETH
  snapshots are up to 0.74\,s later than the BTC snapshots they are joined with in 24.5\%
  of seconds, a small look-ahead that any realistic delay removes.
  \item \textbf{Horizon and trading rule (RQ5).} Longer horizons (30, 60, 300\,s) do not
  improve directional skill. Their occasional profits come from the 18 April sell-off and
  are fragile to the choice of thresholds. The trading rule matters far more than the
  architecture: separating entry from exit turned five of the six 10\,s runs from
  losses of 39--49\% into gains under the backtest's assumptions.
\end{enumerate}

Overall, the LOB contains genuine predictive information about the next mid-price move,
but at one-second resolution that information decays faster than it can be traded after
a 1.5\,bp fee. The most consequential choices in this study were methodological:
out-of-sample selection of trading parameters, walk-forward testing, and modelling
execution delay.

\paragraph{Limitations and future work.} The data covers 12 days in a single month, and
all test days fall in one week that contains one large sell-off. The backtest fills at the
mid-price with a flat fee and full-capital positions, ignoring queue position, adverse
selection, market impact and borrowing costs for short positions. Only the CNN and
LightGBM were evaluated walk-forward. Natural next steps are (i) sub-second or
event-time data, to measure whether the decaying signal can be captured with realistic
latency; (ii) explicit maker/taker execution models with queue position; (iii) using
traded volume (the \texttt{buys}/\texttt{sells} fields) for order-flow and Amihud factors;
and (iv) longer samples covering several market regimes.
